\section{Strain post-processing}
\label{sec:strain-processing}

The Strain window is a separate top-level window that opens
automatically when a Run completes. It computes strain from the
stored displacement field, with its own parameter panel and its own
visualization state (independent of the main window).

You can also open it manually via the \textbf{Open Strain Window}
button in the right sidebar.

\subsection{Strain Parameters panel}

Three editors from top to bottom:

\subsubsection*{Method}

\begin{description}[leftmargin=1.2cm,style=nextline]
  \item[Plane fitting (default)]
    Weighted moving least-squares fit of a first-order polynomial to
    the displacement field inside a virtual strain gauge (VSG).
    Gaussian weights with the VSG radius as $\sigma$. Robust,
    smooths local noise.
  \item[FEM nodal]
    Strain from shape-function derivatives of the FEM mesh. Higher
    spatial resolution but more sensitive to per-node noise.
\end{description}

\subsubsection*{VSG size (plane fitting only)}

Virtual Strain Gauge diameter in pixels. Odd integer. Default 41.
Larger = more smoothing. The field reports radius internally as
$(VSG - 1) / 2$.

\subsubsection*{Strain field smoothing}

Gaussian smoothing of the strain tensor field \emph{after}
computation (not a displacement smoother). Dropdown with four
presets and a warning glyph on the aggressive one:

\begin{center}
\begin{tabular}{lll}
  \toprule
  Preset & $\sigma$ relative to step & Behaviour \\
  \midrule
  Off                  & 0    & No smoothing \\
  Light                & 0.5  & Subtle, preserves fine features \\
  Medium               & 1.0  & Balanced (recommended for noisy data) \\
  Strong $\blacktriangle$ & 2.0 & Aggressive, may blur real gradients \\
  \bottomrule
\end{tabular}
\end{center}

Here ``step'' is the DIC node spacing (the value you set as
\emph{Subset Step} in the main window). Smoothing smaller than
$0.5 \times $ step has essentially no effect because the Gaussian
kernel cannot reach any neighbour node; $2 \times $ step is the
upper bound of ``useful'' smoothing.

\subsubsection*{Strain type}

The strain measure the output fields report:

\begin{description}[leftmargin=1.2cm,style=nextline]
  \item[Infinitesimal (default)]
    $\varepsilon = \frac{1}{2}(\nabla u + \nabla u^T)$. Accurate for
    small deformation only. \emph{Not} zero under finite rigid
    rotation --- it reports $\cos\theta - 1$ as false strain.
  \item[Eulerian--Almansi]
    $e = \frac{1}{2}(I - F^{-T} F^{-1})$. Reported in the current
    (deformed) configuration.
  \item[Green--Lagrangian]
    $E = \frac{1}{2}(F^T F - I)$. Reported in the reference
    configuration. \textbf{Exactly zero under rigid rotation.} Use
    this for any experiment with rotation $>$ 2\textdegree{}.
\end{description}

\subsection{Compute Strain button}

Green primary button below the parameters. Click to recompute strain
with the current panel settings. Runs in a background thread; a
progress bar shows per-frame progress.

If you haven't changed any parameter and the strain is already
computed, the stale warning label stays empty. If you change
anything, a yellow \emph{Params changed --- click Compute Strain}
hint appears.

\subsection{Strain field selector}

Below the parameters, a radio-style selector with all available
fields:

\begin{itemize}
  \item \textbf{Displacement} (shared with main window):
        \texttt{disp\_u}, \texttt{disp\_v}, \texttt{disp\_magnitude},
        \texttt{velocity}
  \item \textbf{Strain components}:
        \texttt{strain\_exx}, \texttt{strain\_eyy}, \texttt{strain\_exy}
  \item \textbf{Derived strains}:
        \texttt{strain\_principal\_max},
        \texttt{strain\_principal\_min},
        \texttt{strain\_maxshear},
        \texttt{strain\_von\_mises},
        \texttt{strain\_rotation} (polar-decomposition rotation
        angle in degrees)
\end{itemize}

Visualization controls (colormap, range, opacity) and physical-unit
settings mirror the main window's but are \emph{independent}: the
strain window remembers its own colormap and range choices per
field.

\subsection{Strain vs main window}

\begin{center}
\begin{tabular}{p{4cm}p{4cm}p{4cm}}
  \toprule
  & Main window & Strain window \\
  \midrule
  Shows & Displacement field & Strain fields (plus disp) \\
  State  & Shared AppState & Private viz state \\
  Export button & Export Results & Export Results (same dialog) \\
  Frame navigator & Shared & Private \\
  \bottomrule
\end{tabular}
\end{center}
